\begin{table}[htbp]
\centering
\caption{Winner-to-Runner-up Bid Gap (18-month window)}
\label{tab:runnerup_gap}
\begin{adjustbox}{max width=\textwidth}
\begin{threeparttable}
\small
\begin{tabular}{lcccc}
\toprule
 & Full sample & $N=2$ firms & $N=3$ firms & $N \geq 5$ firms \\
\midrule
\multicolumn{5}{l}{\textit{Normalized gap: $(p^{(2)} - p^{(1)})/p^\text{ref}$}} \\
$g65 \times Pre$ & -0.0271*** & 0.0071 & -0.0121 & -0.0104* \\
 & (0.0047) & (0.0110) & (0.0097) & (0.0054) \\
\addlinespace
\multicolumn{5}{l}{\textit{Log-gap: $\log(1 + \text{normalized gap})$}} \\
$g65 \times Pre$ & -0.0188*** & 0.0035 & -0.0092 & -0.0082** \\
 & (0.0030) & (0.0068) & (0.0060) & (0.0037) \\
\addlinespace
Observations (normalized gap) & 471,344 & 69,634 & 81,249 & 245,002 \\
\midrule
Item + PBU + Month FE & YES & YES & YES & YES \\
\bottomrule
\end{tabular}
\begin{tablenotes}
\small
\item \textit{Notes:} The \textit{normalized gap} is the difference
between the second-lowest (runner-up) and lowest (winner) phase-2 bids,
normalized by the buyer's reference price. Under more aggressive bidding
the runner-up's bid sits closer to the winner's, and the gap shrinks.
The \textit{log-gap} variant is $\log(1 + \text{gap})$, which is more
robust to right-skewness of the raw gap distribution.
Columns (2)--(4) condition on the exact number of participating firms
to hold the extensive margin approximately fixed. A negative
$\mathit{g65}\times\mathit{Pre}$ coefficient means the runner-up-to-winner
gap narrows under open tenders: direct bidder-level evidence of
intensifying competitive pressure within auctions. Standard errors
clustered at the item level.
*** \textit{p}$<$0.01, ** \textit{p}$<$0.05, * \textit{p}$<$0.1.
\end{tablenotes}
\end{threeparttable}
\end{adjustbox}
\end{table}
